\section{The noise floor}\label{sec:noise}

Appendix Table~\ref{tab:t27} is the supervision problem underneath every routing analysis in this paper: the rows a router would learn from are the contested ones, and those are exactly the rows that flip between two runs of the same condition. On \texttt{cls$\cdot$B0}, rerunning one condition changes the outcome of 10.3--14.3\% of tasks, and 86 of 224 tasks flip in at least one of the six replicated arms. Those flips are not spread evenly. They concentrate in the rows where more than one mode succeeds: 81.8\% of those flip in at least one arm, against 10.3\% of the rest. Both figures are unions over the six replicated arms, so both rise with the number of arms replicated while their ratio falls as they do (17.4$\times$ at two arms, 7.9$\times$ at six); neither the rates nor the enrichment is an arm-count invariant.

\paragraph{What was replicated.} Twenty-five same-condition pairs in six of the eight cells (Appendix Table~\ref{tab:rerun}): all six arms of \texttt{cls$\cdot$B0} ($n=224$), \texttt{red$\cdot$B0} ($n=203$) and \texttt{wa\_red$\cdot$B1} ($n=104$), three arms of \texttt{cls$\cdot$B1}, and two each of \texttt{red$\cdot$B1} and \texttt{cls$\cdot$B2}. \texttt{red$\cdot$B2} and \texttt{wa\_red$\cdot$B0} carry none, and no band is imported into them. How far a rerun moves is a property of the cell, not of the serving path. Re-running one condition moves the set of solved tasks by 4.46--7.59pp in each direction on \texttt{cls$\cdot$B0} (10.3--14.3\% of tasks change outcome), by 1.97--6.90pp on \texttt{red$\cdot$B0} (4.9--11.3\%), and by up to 7.69pp on the locally served \texttt{wa\_red$\cdot$B1} (1.9--9.6\%), but by at most 1.97pp on the local VisualWebArena cells (at most 3.5\%). The threshold below is computed on the \texttt{cls$\cdot$B0} band because that cell carries most of our contrasts; where a cell has its own band, that band should be preferred.

\looseness=-1 \paragraph{Band versus threshold.} Two quantities could each be called a noise floor, and only one of them is one. The \emph{observed band} is how much mean success rate actually moved across those pairs: 0.89--2.68pp. Each of those is \emph{one draw} of a random quantity rather than a bound on it, so reading a 2.2pp effect against a 2.68pp ``measured floor'' compares a draw to a draw. The \emph{resolution threshold} asks a different question: how large must an effect be before a single rerun would be unlikely to produce it unaided? Under the exchangeability null in which each of the $d$ discordant tasks falls either way with probability one half, $\mathrm{SD}(\Delta\mathrm{SR}) = \sqrt{d}/n$ over the $n$ scored tasks, giving 2.14--2.53pp. That is the SD of the \emph{signed} difference, whereas the observed band above is a folded absolute one; the like-for-like null quantity is $\mathbb{E}|\Delta| = \sigma\sqrt{2/\pi}$, or 1.71--2.02pp, still the same order as what repetition actually delivered. An effect has to reach roughly 3.5--4.2pp to clear the threshold. We use the threshold throughout, and quote the observed band only as what repetition happened to deliver. The threshold is $1.645\sigma$, so it is arithmetically a one-sided 5\% bound under that null and we use it as one. What it is not is a test of any hypothesis about representations: the null it rejects is that two runs of one condition were exchangeable, not that two modes are the same. A rerun band states the resolution of this instrument, meaning how small a difference it cannot separate. Where this paper needs an inferential claim it uses a paired interval on the contrast itself (Appendix Tables~\ref{tab:t16} and \ref{tab:t17}). The null behind the threshold models only exchangeability of the two runs, not environment drift; and site state accumulates rather than resets, so drift pushes in one direction and adds to the movement the null already allows. So 3.5--4.2pp is itself a lower bound on what repetition could deliver.

\paragraph{Mean differences versus set differences.} Two runs of one condition also differ as \emph{sets}: $|\{a\text{ solves}\} \setminus \{b\text{ solves}\}|$ runs 4.46--7.59pp on the six \texttt{cls$\cdot$B0} arms. A union ceiling moves by that functional, so an added arm is judged against it in \S\ref{sec:upperbound} rather than against the threshold above. The two must never be compared to each other. A mean-difference threshold applied to a set-difference gain would be arithmetic across estimands, which is how the same replicate can appear to license opposite conclusions.

\looseness=-1 \paragraph{Sources of run-to-run movement.} Sampling is not the answer: every condition in this paper decodes at temperature $0$ with a fixed seed. Inference is nonetheless not bit-reproducible in general \citep{he2025nondeterminism,yuan2025numerical}, and language models are separately sensitive to formatting choices that carry no information \citep{sclar2024promptformat}. Two sources survive a source-level audit of the pipeline. The first is element-id churn. The accessibility-tree payload is keyed by native browser node ids that are reassigned per snapshot; the model is sensitive to those id tokens; and a replay that shuffles ids while holding the page fixed changes which element is chosen on 20.0\% of B1 steps and 12.5\% of B0 steps. Because churn is a property of the text payload and not of the prompt, it reaches exactly the two modes carrying an accessibility tree, DOM and DOM+sprompt. The modes carrying \texttt{[SOM\_MARKS]} are keyed $1\ldots K$ by position and are unaffected. The second applies only to B0. A hosted mixture-of-experts endpoint routes and batches server-side, and no client-side setting controls that. B1 and B2 are served locally and decode greedily, and their steps replay bit-identically (133/133 steps). Their \emph{episodes} still move, because an episode also depends on site state, wall-clock and session lifetime. What every band here names is therefore run-to-run movement including environment drift, never decoding stochasticity. We report no split of the total across these sources. An earlier linear decomposition was withdrawn as a category error, and the per-channel figures above are what each channel moves on its own rather than shares of one budget.

\paragraph{Consequences for the rest of the paper.} Every comparison after it is read against these bands, whether between modes, between deployment classes, or between a representation and a rerun. An arm added for its representation therefore has to be priced against an arm added by repetition (\S\ref{sec:upperbound}), and the tasks a router would be trained on are the least stable rows in the dataset (\S\ref{sec:lowerbound}).
